% Avsnitt 16.1-16.7, ur lectures/L16/appendix/a_can_controller.md.

\section{En nod, två roller}
\label{c16:sec:roles}

En nod antingen \textbf{sänder} sin egen begärda ram eller \textbf{tar emot} någon annans, aldrig
båda för samma ram. En registrerad \code{role} (\code{'1'} = sänder, \code{'0'} = tar emot), låst en
gång per ram på den flank som lämnar \code{STATE_IDLE} för \code{STATE_START} och hållen tills ramen
är slut, avgör vilka av delblockens portar som spelar roll. Båda vägarna ut ur \code{STATE_IDLE}
kräver \textbf{en buss som har legat ledig}: med det etablerat startar \code{tx_req = '1'} en
sändning (\code{role = '1'}), och i annat fall startar ett dominant \code{rx_bus_s2} (någon annans
SOF) en mottagning (\code{role = '0'}). Följ \code{role = '1'}-vägen i det här kapitlet.

\subsection*{Varför ”en buss som har legat ledig” gör verkligt arbete}
Villkoret förtjänar sin plats två gånger om, av olika skäl vid vardera ingången.

\paragraph{Vid mottagaringången}
En dominant bit i sig betyder inte att en ram börjar; de flesta bitarna i de flesta ramar är
dominanta. \code{STATE_IDLE} måste skilja den dominanta bit som \emph{inleder} en ram från de många
dominanta bitarna inuti en, och en nod i \code{STATE_IDLE} har ingen ramkontext att skilja dem åt
med.

Två frestande regler misslyckas båda:
\begin{itemize}
  \item \textbf{En dominant nivå.} Slår till på vilken dominant bit som helst, så en nod som sitter i
    \code{STATE_IDLE} medan en annan nod sänder försöker omedelbart ”ta emot” en ram som redan är
    halvvägs klar.
  \item \textbf{En flank från recessiv till dominant.} Låter som \secref{c10:sec:frameformat}:s
    beskrivning av SOF, men bitstoppning tvingar fram en övergång efter var femte lika bit, så
    flanker i båda riktningarna, den från recessiv till dominant inräknad, förekommer ständigt mitt
    inne i en ram. Flanktestet slår till inom några få bitar från nivåtestet.
\end{itemize}

Regeln som faktiskt fungerar är en \textbf{räkning av recessiva bitperioder}, och dess tröskel kräver
ett steg mer omsorg än stoppbitsregeln ensam. Inuti det stoppade området garanterar
\secref{c10:sec:stuffing}:s regel ingen följd längre än fem lika bitar. Men svansen är ostoppad och
nästan helt recessiv: en kvitterad ram slutar med ACK-avgränsaren plus EOF, åtta recessiva
bitperioder i rad, och en ram som ingen kvitterar visar tio (CRC-avgränsare, recessiv ACK-lucka,
ACK-avgränsare, EOF). En tröskel på sex till tio skulle förklara bussen ledig medan en rams egen svans
fortfarande passerade över den, och ett väntande \code{tx_req} skulle då driva ut SOF över någons
sista EOF-bitar.

Räkna recessiva bitperioder i \code{STATE_IDLE} och kräv \textbf{elva}, samma antal som riktig CAN
ger en nod som måste bedöma om bussen är ledig utan någon ramkontext: den kvitterade svansen på åtta
plus standardens intermission på tre bitar, och förbi även den tio bitar långa svansen hos en
okvitterad ram. Först när räkningen är full betyder ett dominant \code{rx_bus_s2} SOF. Två följder
för resten av designen:
\begin{itemize}
  \item \code{bit_timer} måste fortsätta gå i \code{STATE_IDLE}, eftersom räkningen sker i
    bitperioder.
  \item Reset måste starta räkningen \textbf{full}. En nod som kommer ur reset har ingen anledning
    att tro att en ram pågår, och startade den tom skulle den strunta i den första ramen på bussen.
\end{itemize}

\paragraph{Vid sändaringången}
Samma räkning svarar på en annan fråga: får den här noden börja driva bussen över huvud taget? Riktig
CAN säger att en nod bara får börja sända på en ledig buss, och skälet är inte artighet. Arbitrering
avgör en tvist enbart mellan noder som startade vid \emph{samma} SOF, genom att jämföra vad var och
en sände mot vad bussen läser, bit för bit, ned genom identifieraren. En nod som startar mitt inne i
en ram deltar aldrig i den tävlingen; den driver bara dominanta bitar över någon annans fält. Ramen
som pågår korrumperas, och skadan dyker upp i andra änden som ett stoppbitsbrott eller ett CRC-fel,
utan att något pekar tillbaka på noden som orsakade den. Så \code{tx_req} väntar också på räkningen.

Tillsammans håller de här villkoren \code{error} observerbar. En nod som förlorar arbitreringen
(\chapref{c17:ch}) faller tillbaka till \code{STATE_IDLE} medan vinnaren fortfarande driver bussen,
och \code{STATE_START} nollställer \code{error}. Under någon av de två misslyckade mottagarreglerna
ovan skulle den gå in i \code{STATE_START} igen inom några få bitperioder och sudda den \code{error}
den just höjt; med sändaringången ogrindad skulle en anropare som fortfarande håller \code{tx_req}
hög sudda den redan nästa cykel. Att grinda båda betyder att förloraren väntar, korrekt, på nästa
ram, och att en anropare som pollar \code{error} då och då fortfarande ser felet.

% ------------------------------------------------------------------------------------------
\section{Ramen, fält för fält}
\label{c16:sec:fields}

\lecfig[0.92]{lectures/L16/appendix/images/fsm.png}{\code{can_controller}s
sändningstillståndsmaskin.}{c16:fig:fsm}

Varje \textbf{stoppat} fält går genom \code{tx_shift_reg} som en grupp, sänd MSB först. Varje fält
får \textbf{två} tillstånd: ett \code{STATE_LOAD_*} på en cykel som pulsar \code{tx_shift_reg.load},
och sedan skifttillståndet som väntar på \code{tx_shift_reg.done}.

\begin{booktable}{@{}p{11.6em}p{5em}r>{\raggedright\arraybackslash}p{8.6em}c@{}}
  \thead{Tillstånd} & \thead{Fält} & \thead{Bitar} & \thead{Sänt värde} & \thead{Stoppat} \\
  \midrule
  \code{STATE_START} / \code{STATE_SOF} & SOF & 1 & dominant \code{'0'} & ja \\
  \code{STATE_LOAD_ARB_HI} / \code{STATE_ARB_HI} & ID hög & 8 & \code{tx_id(10 downto 3)} & ja \\
  \code{STATE_LOAD_ARB_LO} / \code{STATE_ARB_LO} & ID låg + RTR & 4 & \code{tx_id(2 downto 0)},
    sedan \code{RTR='0'} & ja \\
  \code{STATE_LOAD_CTRL} / \code{STATE_CTRL} & Kontroll & 6 & \code{IDE='0'}, \code{r0='0'}, sedan
    4 bitars \code{tx_dlc} & ja \\
  \code{STATE_LOAD_DATA} / \code{STATE_DATA} & Databyte & 8 & aktuell byte, upprepad \code{DLC}
    gånger & ja \\
  \code{STATE_CRC_WAIT} & (stabilisering) & 0 & låt \code{crc15} stabilisera sig, lås sedan in den i
    \code{crc_tx} & nej \\
  \code{STATE_LOAD_CRC_HI} / \code{STATE_CRC_HI} & CRC hög & 8 & \code{crc_tx(14 downto 7)} & ja \\
  \code{STATE_LOAD_CRC_LO} / \code{STATE_CRC_LO} & CRC låg & 7 & \code{crc_tx(6 downto 0)} & ja \\
  \code{STATE_CRC_LO_WAIT} & (stabilisering) & 0 & låt \code{crc_valid} stabilisera sig & nej \\
  \code{STATE_CRC_DELIM} & CRC-avgränsare & 1 & recessiv & nej \\
  \code{STATE_ACK_SLOT} & ACK-lucka & 1 & recessiv (mottagaren driver dominant) & nej \\
  \code{STATE_ACK_DELIM} & ACK-avgränsare & 1 & recessiv & nej \\
  \code{STATE_EOF} & EOF & 7 & recessiv & nej \\
\end{booktable}

Noteringar om skiftregistrets laddningar:
\begin{itemize}
  \item \code{tx_shift_reg.data} är en \code{byte_t}, alltid 8 bitar, och sänder de översta
    \code{bit_count} av dem (\chapref{c14:ch}). En grupp som är smalare än en byte är alltså
    \textbf{vänsterjusterad och nollutfylld} vid laddning, och utfyllnaden sänds aldrig. Tabellen
    ovan namnger varje grupp efter dess riktiga bitar; utfyllnaden är mekanisk.
  \item Kontrollbyten laddas som \code{"00" & tx_dlc & "00"} med \code{bit_count = 6}, så de sex
    sända bitarna är de sex översta: \code{0}, \code{0}, \code{DLC[3:0]}.
  \item Den låga CRC-gruppen laddas som \code{crc_tx(6 downto 0) & '0'} med \code{bit_count = 7}.
  \item Den låga ID-gruppen laddas som \code{tx_id(2 downto 0) & '0' & "0000"} med
    \code{bit_count = 4}: de tre återstående bitarna av identifieraren, och sedan den alltid
    dominanta RTR som avslutar arbitreringsfältet.
  \item Data väljer en byte ur \code{tx_data} per \code{data_byte_idx}; loopa
    \code{STATE_LOAD_DATA}/\code{STATE_DATA} en gång per byte, \code{DLC} gånger (hoppa direkt till
    \code{STATE_CRC_WAIT} när \code{DLC = 0}).
  \item \code{tx_dlc} över 8 ligger utanför kontraktet: \secref{c10:sec:frameformat}:s DLC kodar 0
    till 8, och byteloopen går utanför \code{data_byte_idx}s intervall för allt större. Riktiga
    kontrollrar begränsar kodningarna 9 till 15 till åtta byte; den här designen kräver helt enkelt
    att anroparen inte skickar en sådan.
\end{itemize}

De \textbf{ostoppade} avslutande fälten (CRC-avgränsare, ACK-lucka, ACK-avgränsare, EOF) går inte
genom \code{tx_shift_reg}; de är fasta bitantal som drivs direkt, en bit per
\code{bit_timer.bit_done}, identiskt för båda rollerna.

% ------------------------------------------------------------------------------------------
\section{Vad du lägger till i arkitekturen}
\label{c16:sec:declare}

\file{can_controller.vhd} håller redan entiteten (\chapref{c11:ch}), synkroniseraren
\code{meta_prev} (\chapref{c12:ch}) och alla fyra CAN-delblocksinstanserna med de signaler de kopplas
till (\chapref{c12:ch} till \ref{c15:ch}). Den håller inget rambeteende alls. Allt nedan är nytt i
det här kapitlet.

Tillståndstypen och tillståndsmaskinens egna signaler hör hemma i arkitekturens deklarativa del,
ovanför de fyra delblockens signalgrupper:

\begin{vhdlcode}
    -- One state per frame field. Every stuffed field is a LOAD/shift pair; the
    -- trailer fields (CRC delimiter, ACK slot/delim, EOF) are driven directly.
    type state_t is (
        STATE_IDLE,
        STATE_START,
        STATE_SOF,
        STATE_LOAD_ARB_HI, STATE_ARB_HI,
        STATE_LOAD_ARB_LO, STATE_ARB_LO,
        STATE_LOAD_CTRL,   STATE_CTRL,
        STATE_LOAD_DATA,   STATE_DATA,
        STATE_CRC_WAIT,
        STATE_LOAD_CRC_HI, STATE_CRC_HI,
        STATE_LOAD_CRC_LO, STATE_CRC_LO,
        STATE_CRC_LO_WAIT,
        STATE_CRC_DELIM,
        STATE_ACK_SLOT,
        STATE_ACK_DELIM,
        STATE_EOF);

    signal state : state_t;

    -- '1' while transmitting this frame, '0' while only receiving. Latched per frame.
    signal role : std_logic;

    -- Consecutive recessive bit periods seen in STATE_IDLE, saturating at 11. Only once this is
    -- full may the node leave STATE_IDLE at all. Resets to 11, not 0.
    signal idle_bits : natural range 0 to 11;

    -- Data-byte index, and a countdown for the fixed-width trailer fields.
    signal data_byte_idx   : natural range 0 to 8;
    signal field_bit_count : natural range 0 to 7;

    -- Accumulators for the frame being received.
    signal rx_id_acc   : id_t;
    signal rx_dlc_acc  : dlc_t;
    signal rx_data_acc : data_t;

    -- This frame's own CRC, captured at STATE_CRC_WAIT before the engine is fed
    -- the CRC field's own bits and moves on.
    signal crc_tx : crc_t;

    -- The data byte selected by data_byte_idx, MSB-first (byte 0 first).
    signal tx_data_shifted : data_t;
    signal tx_current_byte : byte_t;
\end{vhdlcode}

De tre ackumulatorerna \code{rx_*_acc} används inte förrän i \chapref{c17:ch}; att deklarera dem med
resten håller ihop gruppen i stället för att dela den över två kapitel.

Två konkurrenta tilldelningar plockar ut den databyte som just nu sänds. De hör hemma efter
\code{begin}, vid sidan av delblocksinstanserna:

\begin{vhdlcode}
    tx_data_shifted <= std_logic_vector(shift_left(unsigned(tx_data), 8 * data_byte_idx));
    tx_current_byte <= tx_data_shifted(63 downto 56);
\end{vhdlcode}

\code{shift_left} och \code{unsigned} kommer från \code{numeric_std}, så lägg till
\code{use ieee.numeric_std.all;} i kontextklausulen om den inte redan finns där.

Allt annat i det här kapitlet är de två processer du skriver: en \textbf{kombinatorisk} som driver
delblocken per tillstånd, och en \textbf{sekventiell} som för tillståndsmaskinen framåt och
uppdaterar bokföringen.

\paragraph{Vilken process driver vad}
Uppdelningen är värd att fastställa före genomgången, eftersom varje rad i den landar i en av de två
processerna, och vilken det är ska inte vara någon gåta:
\begin{itemize}
  \item Den \textbf{kombinatoriska} processen driver varje ingång på delblocken (\code{bt_enable},
    \code{bt_resync}, de fyra \code{txsr_*}-ingångarna, de två \code{rxsr_*}-ingångarna i
    \chapref{c17:ch}) och bussparet \code{tx_bus}/\code{bus_en}. Ge dem alla ett förval högst upp och
    skriv över per tillstånd, i samma stil som \code{bit_timer}s egna grenar.
  \item \code{bt_enable} är \code{'1'} i \textbf{varje} tillstånd, \code{STATE_IDLE} inräknat,
    eftersom räkningen i \code{STATE_IDLE} mäts i bitperioder. Porten finns för att \code{bit_timer}
    ska kunna hållas, och den här designen vill aldrig det.
  \item \code{txsr_shift} är \code{bt_bit_done}, men bara under sändning och bara i
    \textbf{skifttillstånden}, aldrig i ett \code{STATE_LOAD_*}-tillstånd. Det är det som håller
    \chapref{c14:ch}:s löfte om att ”\code{load} och \code{shift} utesluter varandra”: ett
    laddtillstånd utfärdar en laddning och inget annat.
  \item Den \textbf{sekventiella} processen äger \code{state} och bokföringen (\code{role},
    \code{idle_bits}, \code{data_byte_idx}, \code{field_bit_count}, \code{crc_tx}, ackumulatorerna
    \code{rx_*_acc}) och varje utgångsport på registersidan: \code{tx_done}, \code{rx_valid},
    \code{error} och \code{rx_*}-registren. Bussparet förblir kombinatoriskt, drivet per tillstånd
    enligt ovan. \code{tx_done} och \code{rx_valid} är pulser på en cykel, så ge dem förvalet
    \code{'0'} vid varje flank och höj dem där de hör hemma.
  \item \code{error} nollställs av \code{reset_s2_n} såväl som av \code{STATE_START}. Den är en nivå
    som en anropare pollar, så den måste läsa \code{'0'} ur reset i stället för vad nu en vippa
    råkade starta upp som.
  \item \code{tx_id}, \code{tx_dlc} och \code{tx_data} låses \textbf{inte} vid \code{STATE_START};
    laddtillstånden läser dem live, vart och ett i det ögonblick dess grupp laddas. Det är en medveten
    förenkling och den lägger skyldigheten på anroparen: håll alla tre stabila från \code{tx_req}
    till \code{tx_done}. Ett registerblock i samma klockdomän gör det gratis, och det är precis den
    anropare \secref{c11:sec:interface} förutsätter.
\end{itemize}

% ------------------------------------------------------------------------------------------
\section{Den kombinatoriska processen}
\label{c16:sec:comb}

Den driver delblockens ingångar och bussparet utifrån \code{state}, med förvalen först och
överskrivningarna per tillstånd efter, och den är kompakt nog att visas hel. Tre saker att se i den
innan koden. Varje \code{STATE_LOAD_*}-arm laddar sin grupp och delar arm med det parade
skifttillståndet, eftersom paret också delar ett \code{bit_count}. \code{txsr_shift} är
\code{bt_bit_done} bara i skifttillstånden, aldrig i ett \code{STATE_LOAD_*}-tillstånd: \code{load}
presenterar redan den nya gruppens första bit under laddcykeln (\chapref{c14:ch}), så ett dedikerat
laddtillstånd på en cykel är det som håller det löftet; att få det här om bakfoten är ramnivåns
version av det enbitsfel i tajmingen som \chapref{c14:ch} varnade för. Och bussen har öppen dränering:
en nod driver bara \emph{aktivt} när biten är dominant, i annat fall släpper den, så att en annan nod
kan dra samma ledning dominant under samma period.

\begin{vhdlcode}
    COMBINATIONAL_PROCESS: process(state, role, bt_bit_done, tx_id, tx_dlc,
                                   tx_current_byte, crc_tx, txsr_tx_bit) is
    begin
        -- Defaults: every sub-block input, every cycle, then per-state overrides.
        bt_enable      <= '1';            -- Never held; the idle count needs it.
        bt_resync      <= '0';
        txsr_load      <= '0';
        txsr_data      <= (others => '0');
        txsr_bit_count <= "0000";
        txsr_shift     <= '0';
        rxsr_enable    <= '0';            -- The receive side arrives in L17.
        rxsr_bit_count <= "0000";
        tx_bus         <= '1';            -- Release the bus unless driving dominant.
        bus_en         <= '0';

        -- Each STATE_LOAD_* arm loads its chunk. The paired shifting state keeps
        -- the same bit_count applied, so grouping the pair in one arm is natural.
        case state is
            when STATE_START =>
                bt_resync <= '1';         -- Adopt SOF as this node's bit boundary.
                if (role = '1') then
                    txsr_load      <= '1';
                    txsr_data      <= "00000000";               -- SOF: one dominant bit.
                    txsr_bit_count <= "0001";
                end if;
            when STATE_LOAD_ARB_HI | STATE_ARB_HI =>
                if ((role = '1') and (state = STATE_LOAD_ARB_HI)) then
                    txsr_load <= '1';
                    txsr_data <= tx_id(10 downto 3);
                end if;
                txsr_bit_count <= "1000";
            when STATE_LOAD_ARB_LO | STATE_ARB_LO =>
                if ((role = '1') and (state = STATE_LOAD_ARB_LO)) then
                    txsr_load <= '1';
                    txsr_data <= tx_id(2 downto 0) & '0' & "0000";  -- ID low, then RTR.
                end if;
                txsr_bit_count <= "0100";
            when STATE_LOAD_CTRL | STATE_CTRL =>
                if ((role = '1') and (state = STATE_LOAD_CTRL)) then
                    txsr_load <= '1';
                    txsr_data <= "00" & tx_dlc & "00";              -- IDE, r0, DLC.
                end if;
                txsr_bit_count <= "0110";
            when STATE_LOAD_DATA | STATE_DATA =>
                if ((role = '1') and (state = STATE_LOAD_DATA)) then
                    txsr_load <= '1';
                    txsr_data <= tx_current_byte;
                end if;
                txsr_bit_count <= "1000";
            when STATE_LOAD_CRC_HI | STATE_CRC_HI =>
                if ((role = '1') and (state = STATE_LOAD_CRC_HI)) then
                    txsr_load <= '1';
                    txsr_data <= crc_tx(14 downto 7);
                end if;
                txsr_bit_count <= "1000";
            when STATE_LOAD_CRC_LO | STATE_CRC_LO =>
                if ((role = '1') and (state = STATE_LOAD_CRC_LO)) then
                    txsr_load <= '1';
                    txsr_data <= crc_tx(6 downto 0) & '0';
                end if;
                txsr_bit_count <= "0111";
            when others =>
                null;
        end case;

        -- Shifting states advance the register once per bit period; a LOAD
        -- state never does, which keeps load and shift mutually exclusive.
        if (role = '1') then
            case state is
                when STATE_SOF | STATE_ARB_HI | STATE_ARB_LO | STATE_CTRL
                   | STATE_DATA | STATE_CRC_HI | STATE_CRC_LO =>
                    txsr_shift <= bt_bit_done;
                when others =>
                    null;
            end case;
        end if;

        -- Driving the bus, open-drain: only the owning states, only dominant
        -- bits. A receiver's one dominant bit is the ACK slot (L17).
        case state is
            when STATE_SOF
               | STATE_LOAD_ARB_HI | STATE_ARB_HI
               | STATE_LOAD_ARB_LO | STATE_ARB_LO
               | STATE_LOAD_CTRL   | STATE_CTRL
               | STATE_LOAD_DATA   | STATE_DATA
               | STATE_LOAD_CRC_HI | STATE_CRC_HI
               | STATE_LOAD_CRC_LO | STATE_CRC_LO =>
                if ((role = '1') and (txsr_tx_bit = '0')) then
                    bus_en <= '1';
                    tx_bus <= '0';
                end if;
            when others =>
                null;
        end case;
    end process;
\end{vhdlcode}

\begin{warning}
Avgränsningen av bussdrivningen i den sista \code{case}-satsen förtjänar sin egen mening:
\code{txsr_tx_bit} är en nivå och håller sitt senaste värde i all evighet, så en oavgränsad version
fortsätter driva efter den grupp som producerade den. En ram vars CRC slutar på en dominant bit
skulle fortfarande dra ned bussen genom hela svansen, och en nod som sitter i \code{STATE_IDLE} med
\code{role} fortfarande \code{'1'} från sin förra ram skulle hålla hela bussen dominant för alltid.
Svansens fält är recessiva för båda rollerna och driver ingenting, med det enda undantaget en
mottagares ACK-lucka (\chapref{c17:ch}).
\end{warning}

% ------------------------------------------------------------------------------------------
\section{Den sekventiella processen}
\label{c16:sec:seq}

Den äger \code{state}, bokföringen och utgångsportarna. Skelett, reset och förval först:

\begin{vhdlcode}
    SEQUENTIAL_PROCESS: process(clock, reset_s2_n) is
    begin
        if (reset_s2_n = '0') then
            state           <= STATE_IDLE;
            role            <= '0';
            idle_bits       <= 11;        -- Full: an idle bus is assumed at power-up.
            data_byte_idx   <= 0;
            field_bit_count <= 0;
            rx_id_acc       <= (others => '0');
            rx_dlc_acc      <= (others => '0');
            rx_data_acc     <= (others => '0');
            crc_tx          <= (others => '0');
            tx_done         <= '0';
            rx_valid        <= '0';
            error           <= '0';
            rx_id           <= (others => '0');
            rx_dlc          <= (others => '0');
            rx_data         <= (others => '0');
        elsif (rising_edge(clock)) then
            tx_done  <= '0';              -- The two pulse outputs default low;
            rx_valid <= '0';              -- error is a level and holds.

            case state is
\end{vhdlcode}

Allt nedan är armar i den här enda \code{case}-satsen. Först de två tillstånd varje ram passerar en
gång:

\begin{vhdlcode}
                when STATE_IDLE =>
                    if (bt_sample = '1') then
                        if (rx_bus_s2 = '1') then
                            if (idle_bits /= 11) then
                                idle_bits <= idle_bits + 1;
                            end if;
                        else
                            idle_bits <= 0;
                        end if;
                    end if;
                    if (idle_bits = 11) then
                        if (tx_req = '1') then
                            role      <= '1';
                            idle_bits <= 0;
                            state     <= STATE_START;
                        elsif (rx_bus_s2 = '0') then
                            role      <= '0';    -- Someone else's SOF: receive (L17).
                            idle_bits <= 0;
                            state     <= STATE_START;
                        end if;
                    end if;
                when STATE_START =>
                    error         <= '0';
                    data_byte_idx <= 0;
                    rx_data_acc   <= (others => '0');
                    state         <= STATE_SOF;
\end{vhdlcode}

\code{STATE_IDLE} räknar recessiva bitperioder vid sampelpunkten och mättar vid 11; ett dominant
sampel nollställer räkningen. Ingenting lämnar tillståndet förrän räkningen är \textbf{full}, vad
\code{tx_req} och \code{rx_bus_s2} än gör, så en \code{tx_req}-puls som anländer medan bussen är
upptagen \textbf{förkastas}, den köas inte; anroparen begär om. \code{STATE_START} behöver ingen kod
för SOF-laddningen eller resynken, eftersom den kombinatoriska processen utfärdar båda utifrån
tillståndet självt; uppsättning som måste \emph{minnas}, att nollställa \code{error}, byteindexet och
ackumulatorn för mottagen data, är det som landar här. \code{rx_data_acc} behöver nollställningen
eftersom en mottagen ram bara skriver de byte dess DLC täcker (\chapref{c17:ch}): en fembytesram som
anländer efter en åttabytesram måste rapportera nollor i byte 5 till 7, inte den förra ramens rester.
\code{rx_id_acc} och \code{rx_dlc_acc} skrivs alltid i sin helhet, så de behöver ingen.
(\code{crc15} nollställs också av det här tillståndet; se \secref{c16:sec:crcclear}.)

Sedan LOAD/skift-paren. Det här är formen varje stoppat fält följer, visad en gång för
identifierarens höga grupp; sändarhalvan av ett LOAD-tillstånd går vidare efter sin enda cykel, och
skifttillståndet väntar på \code{tx_shift_reg.done}.

\begin{aside}
En tajmingnotering innan vågformerna överraskar dig: \chapref{c14:ch} beskrev \code{load} som att den
utfärdas i samma ögonblick som den förra gruppens sista \code{bit_done}, och här utfärdas den två
klockflanker senare, en för att tillståndsmaskinen ska hinna se \code{txsr_done} och en för själva
LOAD-tillståndet. Varje grupps första bit når därför ledningen omkring 40~ns in i sin bitperiod på
1000~ns; sampelpunkten vid 70~\% absorberar det med marginal, och förskjutningen ackumuleras inte,
eftersom varje grupp startar om från samma \code{bit_done}.
\end{aside}

\begin{vhdlcode}
                when STATE_SOF =>
                    if (role = '1') then
                        if (txsr_done = '1') then
                            state <= STATE_LOAD_ARB_HI;
                        end if;
                    end if;                       -- role = '0': L17.
                when STATE_LOAD_ARB_HI =>
                    if (role = '1') then
                        state <= STATE_ARB_HI;    -- One cycle; the load is issued.
                    end if;                       -- role = '0': L17.
                when STATE_ARB_HI =>
                    if (role = '1') then
                        if (txsr_done = '1') then -- Arbitration watching joins here in L17.
                            state <= STATE_LOAD_ARB_LO;
                        end if;
                    end if;                       -- role = '0': L17.
                when STATE_LOAD_ARB_LO =>
                    if (role = '1') then
                        state <= STATE_ARB_LO;
                    end if;                       -- role = '0': L17.
                when STATE_ARB_LO =>
                    if (role = '1') then
                        if (txsr_done = '1') then -- Also watched in L17; RTR is dominant,
                            state <= STATE_LOAD_CTRL;  -- so it can never be the bit that loses.
                        end if;
                    end if;                       -- role = '0': L17.
\end{vhdlcode}

Kontroll- och datatillstånden har samma form med två beslut påhängda: kontrollfältets \code{done}
konsulterar DLC:n, och dataparet loopar en gång per byte:

\begin{vhdlcode}
                when STATE_LOAD_CTRL =>
                    if (role = '1') then
                        state <= STATE_CTRL;
                    end if;                       -- role = '0': L17.
                when STATE_CTRL =>
                    if (role = '1') then
                        if (txsr_done = '1') then
                            if (unsigned(tx_dlc) = 0) then
                                state <= STATE_CRC_WAIT;    -- No data field at all.
                            else
                                state <= STATE_LOAD_DATA;
                            end if;
                        end if;
                    end if;                       -- role = '0': L17.
                when STATE_LOAD_DATA =>
                    if (role = '1') then
                        state <= STATE_DATA;
                    end if;                       -- role = '0': L17.
                when STATE_DATA =>
                    if (role = '1') then
                        if (txsr_done = '1') then
                            if (data_byte_idx + 1 = to_integer(unsigned(tx_dlc))) then
                                state <= STATE_CRC_WAIT;
                            else
                                data_byte_idx <= data_byte_idx + 1;
                                state         <= STATE_LOAD_DATA;
                            end if;
                        end if;
                    end if;                       -- role = '0': L17.
\end{vhdlcode}

Sedan CRC:n. \code{STATE_CRC_WAIT} är en tom cykel så att den sista databiten hinner integreras klart
innan \code{crc_value} läses, och dess verkliga uppgift är låsningen: båda CRC-grupperna laddas från
\code{crc_tx}, aldrig från den levande motorn (\secref{c16:sec:crccapture} förklarar vad som går fel
annars). De två CRC-paren har arbitreringsparens form igen:

\begin{vhdlcode}
                when STATE_CRC_WAIT =>
                    crc_tx <= crc_value;          -- Capture before the engine moves on.
                    state  <= STATE_LOAD_CRC_HI;
                when STATE_LOAD_CRC_HI =>
                    if (role = '1') then
                        state <= STATE_CRC_HI;
                    end if;                       -- role = '0': L17.
                when STATE_CRC_HI =>
                    if (role = '1') then
                        if (txsr_done = '1') then
                            state <= STATE_LOAD_CRC_LO;
                        end if;
                    end if;                       -- role = '0': L17.
                when STATE_LOAD_CRC_LO =>
                    if (role = '1') then
                        state <= STATE_CRC_LO;
                    end if;                       -- role = '0': L17.
                when STATE_CRC_LO =>
                    if (role = '1') then
                        if (txsr_done = '1') then
                            state <= STATE_CRC_LO_WAIT;
                        end if;
                    end if;                       -- role = '0': L17.
\end{vhdlcode}

Till sist svansen, där skiftregistret bugar och går och tillståndsmaskinen själv räknar
\code{bt_bit_done}-pulser. \code{STATE_CRC_LO_WAIT} är en klockcykel för en sändare, \textbf{inte} en
bitperiod: \code{tx_shift_reg.done} pulsar vid den avslutande skiftningen, som själv är ett
\code{bit_done}, så sändaren står redan på en bitgräns och svansen kan börja omedelbart; att vänta en
hel period skulle trycka ut en extra recessiv bit på ledningen. (En mottagare väntar \emph{däremot}
ett helt \code{bit_done} här, och kontrollerar \code{crc_valid} medan den väntar; \chapref{c17:ch}.)

\begin{vhdlcode}
                when STATE_CRC_LO_WAIT =>
                    if (role = '1') then
                        state <= STATE_CRC_DELIM; -- One clock cycle, not a bit period.
                    end if;                       -- role = '0': L17 checks crc_valid here.
                when STATE_CRC_DELIM =>
                    if (bt_bit_done = '1') then
                        state <= STATE_ACK_SLOT;
                    end if;
                when STATE_ACK_SLOT =>
                    if (bt_bit_done = '1') then   -- Transmitter releases; a receiver
                        state <= STATE_ACK_DELIM; -- acknowledges here (L17).
                    end if;
                when STATE_ACK_DELIM =>
                    if (bt_bit_done = '1') then
                        field_bit_count <= 6;     -- Arm the 7-bit EOF countdown.
                        state           <= STATE_EOF;
                    end if;
                when STATE_EOF =>
                    if (bt_bit_done = '1') then
                        if (field_bit_count = 0) then
                            if (role = '1') then
                                tx_done <= '1';
                            end if;                -- role = '0': L17 latches rx_* here.
                            state <= STATE_IDLE;
                        else
                            field_bit_count <= field_bit_count - 1;
                        end if;
                    end if;
            end case;
        end if;
    end process;
\end{vhdlcode}

Svansens tillstånd behöver inget \code{role}-test för att gå vidare, eftersom båda rollerna räknar
samma hela bitperioder genom samma fält; bara det som händer \emph{inuti} dem skiljer sig, och de
skillnaderna är alla \chapref{c17:ch}:s.

% ------------------------------------------------------------------------------------------
\section{Att fånga CRC:n}
\label{c16:sec:crccapture}

\code{STATE_CRC_WAIT} låser \code{crc_value} i \code{crc_tx}, och båda CRC-grupperna laddas från
\code{crc_tx} snarare än från motorn. Den omvägen ser överflödig ut och är det inte, av ett skäl värt
att följa noga, eftersom det är det enda stället i den här designen där två regler som var för sig är
korrekta tillsammans bildar en bugg.

Regel ett: motorn fortsätter gå genom CRC-fältet, eftersom det är det som återför den till noll
(\secref{c16:sec:crcfeed}). Regel två: varje grupp laddas när dess \code{STATE_LOAD_*}-tillstånd
körs.

Sätt nu ihop dem. \code{STATE_LOAD_CRC_HI} körs först och läser rätt värde. Sedan skiftar
\code{STATE_CRC_HI} ut åtta bitar, och var och en av dem går tillbaka in i motorn, som pliktskyldigt
för registret vidare. När \code{STATE_LOAD_CRC_LO} väl körs är \code{crc_value} \textbf{inte längre
ramens CRC}: det är registret åtta bitar längre fram, vilket för en korrekt ram betyder de
återstående sju CRC-bitarna uppskiftade med nollor inskiftade under dem. Ladda
\code{crc_value(6 downto 0)} där, och ramens sista sju CRC-bitar går ut som sju nollor, i varje ram,
alltid.

Felet är tyst på värsta tänkbara sätt. Sändaren är fullkomligt nöjd; ingenting den kan se är fel.
Mottagarens motor når inte noll, så den rapporterar ett CRC-fel, och det uppenbara stället att leta
på är mottagarvägen, vilket är helt i sin ordning. Att låsa värdet en gång, i det ögonblick det är
färdigt och innan något annat rör motorn, tar bort samspelet helt.

% ------------------------------------------------------------------------------------------
\section{Att mata \code{crc15}}
\label{c16:sec:crcfeed}

\code{crc15} måste se varje riktig bit från SOF till och med slutet av \textbf{CRC-fältet}, och inga
andra. Grinda dess enable på \code{tx_shift_reg.bit_valid} (\code{'0'} vid varje grupps avslutande
skiftning, så att den oförändrade biten inte matas in två gånger), och stäng ute stoppbitarna och
mottagarvägen:

\begin{vhdlcode}
crc_enable  <= txsr_bit_valid and not txsr_stuff and role;
crc_data_in <= txsr_tx_bit;
\end{vhdlcode}

Tre saker följer av att det här är en enkel konkurrent tilldelning utan någon tillståndsterm i sig.

\textbf{CRC-fältet matas tillbaka in i motorn, och det är avsiktligt.} Värdet som sänds ligger tryggt
i \code{crc_tx}, så ingenting behöver läsa motorn igen när \code{STATE_CRC_WAIT} väl fångat det, och
uttrycket har ingen tillståndsterm som skulle stoppa matningen vid datafältet. På sändarsidan är
återmatningen bara harmlös: \chapref{c13:ch}:s egenskap generera-och-sedan-kontrollera betyder att ett
meddelande följt av sin egen CRC återför registret till noll, och \code{STATE_START}s nollställning
(\secref{c16:sec:crcclear}) lämnar hur som helst över en fräsch motor till nästa ram. Där
återmatningen förtjänar sin plats är på \textbf{mottagarsidan}: \chapref{c17:ch}:s CRC-kontroll
\emph{är} observationen att registret återvänt till noll när det mottagna CRC-fältet gått in, och den
observationen finns bara därför att CRC-fältet matas tillbaka. En regel för båda rollerna, varje
riktig bit från SOF till och med slutet av CRC-fältet, är det som låter en enda motor generera på en
nod medan den kontrollerar på en annan.

\textbf{Svansen utesluter sig själv.} CRC-avgränsare, ACK-lucka, ACK-avgränsare och EOF går aldrig
genom \code{tx_shift_reg} alls, så \code{txsr_bit_valid} är redan låg över allihop. Ingen
tillståndsterm behövs för att hålla dem utanför checksumman.

\textbf{\code{role} är platshållaren för mottagarsidan.} Med \code{role = '0'} håller det här
uttrycket \code{crc15} avstängd, vilket är korrekt för det här kapitlet (bara sändning) men inte för
den färdiga kontrollern. \Chapref{c17:ch} ersätter hela tilldelningen med en som matar
\code{rx_shift_reg.real_bit} på \code{rx_shift_reg.real_bit_valid} under mottagning; sändarhalvan
ovan är oförändrad av det.

% ------------------------------------------------------------------------------------------
\section{Att nollställa \code{crc15}}
\label{c16:sec:crcclear}

Ännu en konkurrent tilldelning, och det är den som gör motorn säker att återanvända ram efter ram:

\begin{vhdlcode}
crc_clear <= '1' when state = STATE_START else '0';
\end{vhdlcode}

På den lyckliga vägen gör den ingenting alls, eftersom en ram som körs hela vägen redan har återfört
registret till noll. Den finns där för de ramar som \emph{inte} blir klara. \Chapref{c17:ch} ger den
här noden tre sätt att överge en ram halvvägs (den förlorar arbitreringen, den ser ett
stoppbitsbrott, dess CRC-kontroll misslyckas), och vart och ett lämnar motorn med ett halvfärdigt
värde och ingen utsikt att någonsin nå noll av sig själv. Varje ram noden sänder därefter skulle bära
en CRC beräknad ur den resten, och varje ram den tog emot skulle underkännas i kontrollen, tills
någon tryckte på reset.

Två saker är värda att lägga märke till. \code{STATE_START} är rätt hem för den eftersom det redan är
tillståndet ”sätt upp den här ramen”, vid sidan av att nollställa \code{error} och byteindexet;
nollställningen hör hemma där snarare än bland CRC-tillstånden. Och \code{can_controller_tb} kan inte
fånga en utebliven nollställning, eftersom ingen av dess två scenarier låter en nod sända igen efter
ett avbrott. Det gör det här till en regel att få rätt genom resonemang snarare än genom att vänta på
en röd testbänk, samma lärdom som \code{crc_tx} gav i \secref{c16:sec:crccapture}: de tysta felen är
de man ska resonera sig fram till i förväg.
